\documentclass[11pt,a4paper]{article}
\usepackage[utf8]{inputenc}
\usepackage[margin=1in]{geometry}
\usepackage{amsmath}
\usepackage{amsfonts}
\usepackage{amssymb}
\usepackage{graphicx}
\usepackage{xcolor}
\usepackage{fancyhdr}
\usepackage{hyperref}
\usepackage{tcolorbox}
\usepackage{enumitem}

% Colors
\definecolor{primaryblue}{RGB}{31,119,180}
\definecolor{warningred}{RGB}{255,107,107}
\definecolor{successgreen}{RGB}{46,204,113}
\definecolor{infoblue}{RGB}{52,152,219}

% Header/Footer
\pagestyle{fancy}
\fancyhf{}
\rhead{Y=R+S+N Framework Guide}
\lhead{TidyLLM Whitepapers}
\cfoot{\thepage}

\title{\textbf{Understanding the Y=R+S+N Framework}\\
\large{A Simple Guide to Mathematical Content Analysis}}
\author{Rudy Martin \\ TidyLLM Research Project}
\date{\today}

\begin{document}

\maketitle

\begin{tcolorbox}[colback=infoblue!10,colframe=infoblue,title=What You'll Learn]
This guide explains the Y=R+S+N framework in simple terms, so anyone can understand how to evaluate the quality and usefulness of research papers and content.
\end{tcolorbox}

\tableofcontents
\newpage

\section{What is Y=R+S+N?}

\subsection{The Big Picture}
Imagine you're reading a research paper. Some parts are super useful, some parts are just okay, and some parts are confusing or wrong. The Y=R+S+N framework helps you measure how much of each type of content you're dealing with.

\begin{tcolorbox}[colback=successgreen!10,colframe=successgreen,title=Simple Definition]
\textbf{Y=R+S+N} is like a recipe that breaks down any piece of content into three ingredients:
\begin{itemize}
\item \textbf{R} = Really good stuff (Relevant)
\item \textbf{S} = So-so stuff (Superfluous) 
\item \textbf{N} = Not good stuff (Noise)
\end{itemize}
\end{tcolorbox}

\subsection{The Math (Don't Worry, It's Simple!)}

The basic formula is:
\[ Y = R + S + N \]

But we also calculate a "usefulness score":
\[ \text{Y Score} = R + (0.5 \times S) \]

\textbf{Why multiply S by 0.5?} Because the "so-so stuff" is only half as valuable as the "really good stuff."

\section{Breaking Down Each Component}

\subsection{R - Relevant Content (The Good Stuff)}

\begin{tcolorbox}[colback=successgreen!10,colframe=successgreen,title=What is Relevant Content?]
This is content that directly helps you understand the topic. Think of it like the main ingredients in your favorite recipe.
\end{tcolorbox}

\textbf{Examples of Relevant Content:}
\begin{itemize}[leftmargin=*]
\item Clear mathematical formulas that work
\item Step-by-step explanations of methods  
\item Concrete examples with numbers
\item Proven theories and validated results
\item Direct answers to your research questions
\end{itemize}

\textbf{Scoring:} Relevant content typically ranges from 25\% to 85\% of a paper.
\begin{itemize}
\item \textbf{85\%+ R:} Excellent paper, use as primary source
\item \textbf{60-85\% R:} Good paper, useful for research
\item \textbf{40-60\% R:} Okay paper, some useful parts
\item \textbf{Below 40\% R:} Poor paper, be very careful
\end{itemize}

\subsection{S - Superfluous Content (The Okay Stuff)}

\begin{tcolorbox}[colback=orange!10,colframe=orange,title=What is Superfluous Content?]
This is content that's somewhat related but not essential. Like garnish on a dish - nice to have but not the main point.
\end{tcolorbox}

\textbf{Examples of Superfluous Content:}
\begin{itemize}[leftmargin=*]
\item Long introductions explaining basic concepts
\item Extensive literature reviews
\item Historical background information
\item "Future work" sections
\item Acknowledgments and administrative text
\item General discussions without specific conclusions
\end{itemize}

\textbf{Scoring:} Superfluous content typically ranges from 10\% to 45\%.
\begin{itemize}
\item \textbf{10-25\% S:} Focused paper, gets to the point
\item \textbf{25-35\% S:} Normal amount of background
\item \textbf{35-45\% S:} Lots of background, may be unfocused
\end{itemize}

\subsection{N - Noise Content (The Bad Stuff)}

\begin{tcolorbox}[colback=warningred!10,colframe=warningred,title=What is Noise Content?]
This is content that's confusing, wrong, or irrelevant. Like finding a hair in your soup - it doesn't belong there.
\end{tcolorbox}

\textbf{Examples of Noise Content:}
\begin{itemize}[leftmargin=*]
\item Mathematical errors or typos
\item Unclear or ambiguous explanations
\item Contradictory statements
\item Formatting problems that make text unreadable
\item Completely unrelated information
\item Broken equations or missing steps
\end{itemize}

\textbf{Scoring:} Noise content should be as low as possible.
\begin{itemize}
\item \textbf{0-10\% N:} Very clean paper, high quality
\item \textbf{10-20\% N:} Some issues but still usable
\item \textbf{20-35\% N:} Significant problems, use with caution
\item \textbf{Above 35\% N:} Avoid this source
\end{itemize}

\section{Real-World Examples}

\subsection{Example 1: High-Quality Math Paper}
\textbf{Content:} "We present a robust mathematical framework for signal decomposition with comprehensive validation across multiple datasets."

\textbf{Analysis:}
\begin{itemize}
\item \textbf{R = 78\%} (robust framework, mathematical, comprehensive validation)
\item \textbf{S = 15\%} (some background context)  
\item \textbf{N = 7\%} (minimal errors)
\item \textbf{Y Score = 0.86} (Excellent!)
\end{itemize}

\subsection{Example 2: Survey Paper}
\textbf{Content:} "This comprehensive survey provides extensive background information and literature review of information theory applications."

\textbf{Analysis:}
\begin{itemize}
\item \textbf{R = 35\%} (some relevant mathematical concepts)
\item \textbf{S = 55\%} (lots of background and literature review)
\item \textbf{N = 10\%} (minimal errors)
\item \textbf{Y Score = 0.63} (Decent for background reading)
\end{itemize}

\subsection{Example 3: Problematic Paper}
\textbf{Content:} "Complex multidimensional analysis with unclear mathematical formulations and ambiguous theoretical frameworks. Some equations contain formatting errors."

\textbf{Analysis:}
\begin{itemize}
\item \textbf{R = 57\%} (has mathematical content but unclear)
\item \textbf{S = 24\%} (some background)
\item \textbf{N = 19\%} (unclear formulations, formatting errors)
\item \textbf{Y Score = 0.69} (Use with caution!)
\end{itemize}

\section{Context Collapse Risk}

\begin{tcolorbox}[colback=warningred!10,colframe=warningred,title=What is Context Collapse?]
Context collapse happens when bad or confusing information starts to overwhelm the good information, making it hard to understand what's actually true or useful.
\end{tcolorbox}

\textbf{The Risk Formula:}
\[ \text{Context Collapse Risk} = S + (1.5 \times N) \]

\textbf{Why do we multiply N by 1.5?} Because noise is more dangerous than superfluous content - errors hurt more than extra background info.

\subsection{Risk Levels}
\begin{itemize}
\item \textbf{Low Risk (0.0-0.3):} Safe to use, reliable content
\item \textbf{Medium Risk (0.3-0.6):} Be careful, validate key claims
\item \textbf{High Risk (0.6+):} Dangerous, avoid or use extreme caution
\end{itemize}

\subsection{Types of Context Collapse}
\begin{enumerate}
\item \textbf{Context Poisoning:} Misinformation gets mixed in with good info
\item \textbf{Context Distraction:} Too much information overwhelms the main point
\item \textbf{Context Confusion:} Unclear or contradictory explanations
\item \textbf{Context Clash:} Different sources say completely different things
\end{enumerate}

\section{How to Use This Framework}

\subsection{For Students}
\begin{itemize}
\item Use papers with \textbf{Y Score > 0.8} as your primary sources
\item Papers with \textbf{Y Score 0.6-0.8} are good for supporting information
\item Avoid papers with \textbf{Y Score < 0.5} unless you have no other choice
\item Always check the \textbf{Context Collapse Risk} - if it's "High," get a second opinion
\end{itemize}

\subsection{For Researchers}
\begin{itemize}
\item Aim for \textbf{R > 70\%} and \textbf{N < 15\%} in your own papers
\item When citing others, prefer sources with low context collapse risk
\item Use this framework to quickly screen large numbers of papers
\item Document Y scores in your literature review process
\end{itemize}

\subsection{For Business Professionals}
\begin{itemize}
\item When evaluating technical reports, look for high R scores
\item Be suspicious of reports with high N scores (lots of errors/unclear parts)
\item Use this framework to assess the quality of data and recommendations
\item Ask vendors to provide Y=R+S+N analysis of their technical documentation
\end{itemize}

\section{Quick Reference Guide}

\begin{table}[h]
\centering
\begin{tabular}{|l|c|c|l|}
\hline
\textbf{Y Score} & \textbf{Quality} & \textbf{Risk} & \textbf{Action} \\
\hline
0.9+ & Excellent & Low & Use as primary source \\
0.8-0.9 & Very Good & Low & Great supporting source \\
0.7-0.8 & Good & Medium & Useful with validation \\
0.6-0.7 & Fair & Medium & Use carefully \\
0.5-0.6 & Poor & Medium-High & Secondary use only \\
Below 0.5 & Very Poor & High & Avoid if possible \\
\hline
\end{tabular}
\caption{Y Score Interpretation Guide}
\end{table}

\section{Conclusion}

The Y=R+S+N framework is like having a quality detector for research content. It helps you quickly identify:

\begin{itemize}
\item What's worth your time (high R, low N)
\item What needs extra checking (medium scores)
\item What to avoid (high N, high context collapse risk)
\end{itemize}

\begin{tcolorbox}[colback=successgreen!10,colframe=successgreen,title=Remember]
The goal isn't to get perfect scores - it's to make informed decisions about what information to trust and how carefully to verify it.
\end{tcolorbox}

\vfill
\begin{center}
\textit{This framework is part of the TidyLLM research project.\\
For more information, visit the TidyLLM Whitepapers demo.}
\end{center}

\end{document}